\documentclass[11pt,a4paper]{article}
\usepackage[margin=2.2cm]{geometry}
\usepackage[T1]{fontenc}
\usepackage{lmodern}
\usepackage{amsmath,amssymb}
\usepackage{booktabs,tabularx}
\usepackage{enumitem}
\usepackage[most]{tcolorbox}
\usepackage{xcolor}
\usepackage[hidelinks]{hyperref}

\setlist{itemsep=1pt,topsep=3pt}
\newcommand{\W}{\textsf{\textbf{[Wokwi]}}}
\newcommand{\J}{\textsf{\textbf{[JSLinux]}}}
\newcommand{\Pap}{\textsf{\textbf{[Paper]}}}
\newcommand{\ecu}{\emph{The ECU does not see the engine. It sees only signals.}}

\newtcolorbox{sheetinfo}{colback=gray!6,colframe=gray!50,boxrule=0.4pt,
  left=4pt,right=4pt,top=3pt,bottom=3pt,fontupper=\small}

% \exercise{id}{title}{platform}{description}
\newcommand{\exercise}[4]{%
  \par\smallskip\noindent\textbf{#1\quad #2}\hfill#3\par
  \nopagebreak\begin{itemize}[leftmargin=1.2em]#4\end{itemize}}

\title{Digital Signal Processing -- Exercise Outline\\
\large Engine signals as the guiding example, run on Wokwi and JSLinux}
\author{}
\date{Draft, \today}

\begin{document}
\maketitle

\section*{Concept}

Every exercise sheet follows the leitmotif of the lecture: \ecu{}
The exercises therefore build, chapter by chapter, a small \emph{virtual
engine} (signal generators for crankshaft, camshaft, knock, ion current and
lambda sensors) and a small \emph{virtual ECU} that extracts information
from these signals. Sheet 12 (knock detection) assembles all parts.

\subsection*{Platforms}
\begin{description}[leftmargin=1.5em]
  \item[\W] Browser-based microcontroller simulator
    (\url{https://wokwi.com}). Targets: Arduino Uno (8-bit, no FPU, 10-bit
    ADC -- for fixed-point and resource-limit exercises) and ESP32 (FPU,
    12-bit ADC -- for floating-point real-time filtering). Inputs via
    potentiometers, buttons and \emph{Wokwi custom chips} (C API) that act
    as sensor signal generators; outputs via Serial Monitor (CSV), LEDs and
    the logic analyzer (VCD export).
  \item[\J] x86 Linux in the browser (\url{https://bellard.org/jslinux}).
    Plain C99 with \texttt{libm} (\texttt{gcc} or \texttt{tcc}), offline
    analysis of recorded or synthetic data, text/ASCII plots and CSV output.
    Data transfer from Wokwi: copy the Serial Monitor CSV, upload into
    JSLinux.
  \item[\Pap] Pen-and-paper part (derivations, sketches, short answers),
    used to prepare or check the programming part.
\end{description}

\subsection*{Course infrastructure (to be provided once)}
\begin{itemize}
  \item \texttt{engine\_signals.h/.c}: portable C models of all sensor
    signals (60-2 crank wheel, cam, knock = damped resonances + noise,
    ion current, switching lambda sensor), shared by Wokwi and JSLinux.
  \item Wokwi custom chips \texttt{crank-60-2}, \texttt{knock-sensor}, \texttt{lambda-probe}
    built from the same models, with RPM/load/knock-intensity controls.
  \item \texttt{asciiplot.h}: dependency-free terminal plot for JSLinux
    (time signal, spectrum, pole-zero map).
  \item \texttt{csvio.h}: read/write of sample files; a set of
    pre-recorded data files (crank, knock with/without knock, lambda).
  \item Sheet 0 (\emph{Getting started}, 45\,min): blink + serial CSV and
    the engine simulator chip on Wokwi; compile/run/upload/download on
    JSLinux as an optional bonus exercise. Sheet 2 repeats the JSLinux
    setup in a short section for everybody.
\end{itemize}

\subsection*{Sheet format}
Each sheet: learning goals referencing the script, 1--2 paper exercises,
2--3 lab exercises, deliverables (code, plot, short written answers),
optional challenge. Solutions in the same \texttt{.tex} file, switched by
\verb|\ifsolutions|. Expected workload per sheet: 90\,min tutorial
+ about 3\,h self-study.

\subsection*{Overview}
\begin{center}
\small
\begin{tabularx}{\textwidth}{@{}clXcc@{}}
\toprule
\# & Chapter & Lab core & \W & \J\\
\midrule
1 & ECU and Signals & 60-2 crank signal, RPM from tooth timing (bonus: signal zoo) & x & (x)\\
2 & Systems & Black-box property tests, lambda loop with delay & x & x\\
3 & LTI Systems & Convolution, impulse/step response, moving-average RPM & x & x\\
4 & Sampling & ADC sampling, quantisation SNR, time vs.\ angle domain & x & x\\
5 & Nyquist & Aliasing of knock frequencies, anti-aliasing filter & x & x\\
6 & Fourier & Fourier series of crank signal, engine orders & & x\\
7 & Laplace & Knock sensor as resonator, RC filter, step responses & & x\\
8 & DFT/FFT & Own radix-2 FFT, windowing, leakage, cost & x & x\\
9 & Z-Transform & Pole/zero, stability, digital resonator & x & x\\
10 & Digital Filters & FIR/IIR knock band-pass, fixed point vs.\ float & x & x\\
11 & Maximum Entropy & Burg/AR spectrum vs.\ FFT on short windows & & x\\
12 & Knock Detection & Angle-synchronous real-time knock detector & x & x\\
\bottomrule
\end{tabularx}
\end{center}
(x) = optional bonus exercise only. Sheets 0 and 1 (the first lab days) are
Wokwi-only; JSLinux becomes mandatory from Sheet~2 on.

%=====================================================================
\section{Sheet 1 -- The ECU Sees Only Signals}
\begin{sheetinfo}
Goals: describe sensor signals mathematically, classify them, extract a
first piece of information (engine speed) from a raw signal.
\end{sheetinfo}

\exercise{1.1}{Signal classification}{\Pap}{
  \item Classify the five engine signals (continuous/discrete,
    deterministic/stochastic, periodic/aperiodic, analog/digital) and
    justify each entry of Table~1.2; discuss borderline cases
    (\emph{quasi-periodic} knock, crank signal during acceleration).
  \item Relation crank angle vs.\ cam angle (720$^\circ$ vs.\ 360$^\circ$):
    period of each signal at 3000\,rpm.}

\exercise{1.2}{60-2 crank wheel and engine speed}{\W}{
  \item Arduino Uno + \texttt{crank-60-2} chip: capture rising edges by
    interrupt, store tooth periods, print CSV.
  \item Detect the missing-tooth gap (period ratio $>$ 2), count teeth,
    compute crank angle and RPM; light an LED at TDC.
  \item Vary RPM with a potentiometer; observe jitter in the RPM estimate.}

\exercise{1.3}{Cam/crank synchronisation (challenge)}{\W}{
  \item Add the cam signal and determine the 720$^\circ$ cycle phase
    (cylinder identification).}

\exercise{1.B1}{Signal zoo (bonus, optional)}{\J}{
  \item Generate and plot all five sensor models with
    \texttt{engine\_signals.h}; compute mean, RMS, min/max.
  \item Determine numerically whether a signal is periodic (autocorrelation
    peak) -- compare crank at constant speed vs.\ knock signal.}

%=====================================================================
\section{Sheet 2 -- Systems}
\begin{sheetinfo}
Goals: the operator view $y=S\{x\}$; static/dynamic, causal/non-causal,
memoryless/with memory; first contact with linearity and time invariance.
First sheet with mandatory JSLinux exercises: starts with a short JSLinux
setup section (for students who skipped the bonus exercises of Sheets 0--1).
\end{sheetinfo}

\exercise{2.1}{Properties by hand}{\Pap}{
  \item For about 8 systems (e.g.\ $y=2x$, $y=x^2$, $y[n]=x[n+1]-x[n]$,
    three-point moving average, $y(t)=x(t-T_d)$, rectifier, integrator,
    $y[n]=n\,x[n]$) decide static/dynamic, causal, memory, linear,
    time-invariant.}

\exercise{2.2}{Black-box system identification}{\J}{
  \item Provided object file with five unknown systems
    \texttt{S1..S5(x, y, N)}. Students write test harnesses: superposition
    test, shifted-input test, causality test (impulse at $n_0$), memory
    test.
  \item Report: property table with numerical evidence.}

\exercise{2.3}{Engine as a dynamic system}{\J}{
  \item Simulate crankshaft speed as first-order system (inertia) driven
    by fuel quantity; step in fuel $\rightarrow$ gradual RPM change.
  \item Lambda path: fuel $\rightarrow$ engine $\rightarrow$ transport
    delay $\rightarrow$ switching lambda sensor (non-linear, static).}

\exercise{2.4}{Closed lambda loop}{\W}{
  \item ESP32 + \texttt{lambda-probe} chip: implement a two-point
    (bang-bang) controller with integrator; observe the limit cycle of
    Fig.~1.3 in the Serial Monitor; vary the transport delay and explain
    the change of the oscillation period.}

%=====================================================================
\section{Sheet 3 -- LTI Systems}
\begin{sheetinfo}
Goals: impulse response, convolution, step response, difference
equations; LTI systems as the workhorse of DSP.
\end{sheetinfo}

\exercise{3.1}{Convolution by hand}{\Pap}{
  \item Graphical and tabular convolution of short sequences;
    impulse/step response of moving average and of
    $y[n]=a\,y[n-1]+(1-a)\,x[n]$.}

\exercise{3.2}{Convolution engine}{\J}{
  \item Implement \texttt{conv()} and verify commutativity, associativity
    and distributivity numerically; compare with recursive implementation
    of the exponential smoother.
  \item Cascading: two moving averages = triangular impulse response.}

\exercise{3.3}{Smoothing the RPM estimate}{\W}{
  \item Extend Ex.~1.3: moving average over $N$ teeth vs.\ exponential
    smoother; trade-off noise vs.\ delay during an RPM ramp; measure the
    lag in teeth.
  \item Why must the gap tooth be treated separately (violation of time
    invariance in the tooth domain)?}

%=====================================================================
\section{Sheet 4 -- Sampling}
\begin{sheetinfo}
Goals: sampling and quantisation, ADC resolution, SNR, timing jitter,
time-based vs.\ crank-angle-based sampling.
\end{sheetinfo}

\exercise{4.1}{Quantisation}{\Pap}{
  \item Derive $\mathrm{SNR}\approx 6.02N+1.76$\,dB; required ADC bits for
    the knock sensor dynamic range; LSB size for 10-/12-bit ADCs.}

\exercise{4.2}{ADC in practice}{\W}{
  \item Uno/ESP32: timer-driven sampling of a potentiometer and of the
    \texttt{knock-sensor} chip; verify the sampling rate with the logic
    analyzer (toggle pin per sample).
  \item Note: Wokwi runs in simulated time; timing is measured in
    simulated time, not wall-clock time.}

\exercise{4.3}{Quantisation noise and jitter}{\J}{
  \item Quantise a sine with $N=4\ldots16$ bits; measure SNR, compare with
    the formula; add sampling jitter and observe the SNR drop vs.\
    frequency.}

\exercise{4.4}{Time domain vs.\ angle domain}{\J}{
  \item Resample a time-sampled signal during an RPM ramp onto a
    crank-angle grid (linear interpolation using crank edges); show that
    crank-synchronous components become stationary in the angle domain.}

%=====================================================================
\section{Sheet 5 -- Nyquist}
\begin{sheetinfo}
Goals: sampling theorem, aliasing, reconstruction, anti-aliasing filters.
\end{sheetinfo}

\exercise{5.1}{Aliasing by hand}{\Pap}{
  \item Alias frequency of knock components (6--15\,kHz) for several
    sampling rates; minimum $f_s$; folding diagram.}

\exercise{5.2}{Aliasing made visible}{\W}{
  \item Uno ADC (max.\ about 9\,kHz achievable) samples the knock chip at
    6, 8, 11\,kHz resonance: observe the apparent frequency in the CSV;
    switch to ESP32 with higher $f_s$.}

\exercise{5.3}{Reconstruction and anti-aliasing}{\J}{
  \item Sinc interpolation vs.\ zero-order hold vs.\ linear interpolation;
    reconstruction error.
  \item Simulate an RC (and 2nd-order) anti-aliasing filter before
    sampling; residual alias power vs.\ filter order.}

%=====================================================================
\section{Sheet 6 -- Fourier}
\begin{sheetinfo}
Goals: Fourier series, Fourier transform, DTFT; engine orders.
\end{sheetinfo}

\exercise{6.1}{Fourier series of engine signals}{\Pap}{
  \item Fourier coefficients of a rectangular tooth signal; effect of the
    missing teeth on the spectrum (once-per-revolution modulation).
  \item Damped sinusoid (knock) $\rightarrow$ Lorentzian spectrum;
    bandwidth vs.\ damping.}

\exercise{6.2}{Numerical Fourier series and DTFT}{\J}{
  \item Compute Fourier coefficients numerically; Gibbs phenomenon for
    partial sums; DTFT evaluated on a fine grid (no FFT yet).}

\exercise{6.3}{Engine orders}{\J}{
  \item Signal with 2nd and 4th engine order and a fixed structural
    resonance: spectra at 1500/3000/6000\,rpm; distinguish
    speed-proportional orders from fixed resonances.}

%=====================================================================
\section{Sheet 7 -- Laplace}
\begin{sheetinfo}
Goals: transfer functions, poles and zeros, stability and step response
of continuous-time systems; bridge to discretisation.
\end{sheetinfo}

\exercise{7.1}{Sensor models in the $s$-domain}{\Pap}{
  \item Knock sensor as 2nd-order resonator $H(s)$ (resonance frequency,
    damping, Q); RC low-pass; lambda sensor as first-order lag + delay
    ($e^{-sT_d}$, Pad\'e approximation).
  \item Poles, step response, Bode sketch.}

\exercise{7.2}{Numerical verification}{\J}{
  \item Simulate the ODEs (Euler, RK4) and compare with analytic step and
    impulse responses; evaluate $|H(j\omega)|$ and phase on a grid.
  \item Discretise by forward Euler, impulse invariance and bilinear
    transform; compare frequency responses (preview of Sheets 9/10).}

%=====================================================================
\section{Sheet 8 -- DFT and FFT}
\begin{sheetinfo}
Goals: DFT as sampled DTFT, frequency resolution, leakage, windows,
FFT algorithm and complexity.
\end{sheetinfo}

\exercise{8.1}{DFT properties}{\Pap}{
  \item 4- and 8-point DFT by hand; resolution $\Delta f=f_s/N$; zero
    padding vs.\ true resolution; window main lobe/side lobe.}

\exercise{8.2}{Own FFT}{\J}{
  \item Implement direct DFT and iterative radix-2 FFT; verify against each
    other; measure runtime vs.\ $N$ and fit $N^2$ vs.\ $N\log N$.
  \item Rectangular, Hann, Hamming, Blackman windows: leakage of a
    non-bin-centred knock frequency.}

\exercise{8.3}{FFT on the microcontroller}{\W}{
  \item ESP32: 256-point FFT of the knock chip signal in a crank-angle
    window, print the spectrum, identify the knock resonance.
  \item Uno: memory limit -- largest feasible $N$; operation count instead
    of execution time (Wokwi is not cycle-accurate).}

%=====================================================================
\section{Sheet 9 -- Z-Transform}
\begin{sheetinfo}
Goals: $z$-transform, ROC, transfer function $H(z)$, poles/zeros,
stability, frequency response from the pole-zero map.
\end{sheetinfo}

\exercise{9.1}{Analytical}{\Pap}{
  \item Transforms of standard sequences; inverse by partial fractions;
    $H(z)$ of the exponential smoother and of the moving average;
    stability criterion.}

\exercise{9.2}{Pole-zero explorer}{\J}{
  \item Program that reads $H(z)$ coefficients, computes poles/zeros
    (quadratic formula/Durand--Kerner), ASCII pole-zero plot, magnitude
    and phase response, impulse response.}

\exercise{9.3}{Digital resonator for knock frequency}{\W}{
  \item Two-pole resonator at the knock frequency; pole radius $r$
    vs.\ bandwidth; Uno with float vs.\ Q15 fixed point: observe
    instability/limit cycles for $r\to1$.}

%=====================================================================
\section{Sheet 10 -- Digital Filters}
\begin{sheetinfo}
Goals: FIR (window method) and IIR (bilinear transform) design,
structures (direct form, biquad cascade), coefficient quantisation.
\end{sheetinfo}

\exercise{10.1}{Design by hand}{\Pap}{
  \item Specification of a knock band-pass (e.g.\ 6--8\,kHz, $f_s=40$\,kHz);
    FIR order estimate; 2nd-order IIR band-pass via bilinear transform
    with pre-warping.}

\exercise{10.2}{Filter design tool}{\J}{
  \item Implement FIR window design and IIR biquad design; verify with the
    Sheet 9 explorer; export coefficients as C header (float and Q15).
  \item Coefficient quantisation: pole shift and response change for
    8/12/16-bit coefficients.}

\exercise{10.3}{Real-time filtering}{\W}{
  \item ESP32: knock band-pass (FIR and biquad) on the knock chip signal,
    sample-by-sample in the timer ISR; compare outputs.
  \item Uno: Q15 biquad; overflow handling and scaling.}

%=====================================================================
\section{Sheet 11 -- Maximum Entropy}
\begin{sheetinfo}
Goals: autoregressive (AR) modelling, Yule--Walker, Levinson--Durbin,
Burg method; high-resolution spectra from short records.
\end{sheetinfo}

\exercise{11.1}{AR models}{\Pap}{
  \item Yule--Walker equations for AR(1), AR(2); relation of AR(2) poles to
    a resonance; maximum-entropy principle.}

\exercise{11.2}{MEM vs.\ FFT on short windows}{\J}{
  \item Implement Levinson--Durbin and Burg; spectra of 32/64/128-sample
    knock windows vs.\ windowed periodogram.
  \item Model order selection (FPE/AIC); spurious peaks and line splitting
    for too high orders.}

\exercise{11.3}{Resolving two resonances}{\J}{
  \item Two knock modes close in frequency: minimum record length for
    resolution with FFT vs.\ Burg; influence of SNR.}

%=====================================================================
\section{Sheet 12 -- Knock Detection (capstone)}
\begin{sheetinfo}
Goals: combine all previous parts into a real-time, angle-synchronous
knock detector; evaluate it statistically.
\end{sheetinfo}

\exercise{12.1}{Detector design}{\Pap}{
  \item Signal chain: crank sync $\rightarrow$ measurement window
    (e.g.\ 10--70$^\circ$ after TDC) $\rightarrow$ band-pass
    $\rightarrow$ rectification/squaring $\rightarrow$ integration
    $\rightarrow$ normalisation by background level $\rightarrow$
    threshold. Justify each block with the respective chapter.}

\exercise{12.2}{Offline evaluation}{\J}{
  \item Run the detector on recorded data sets with known knock labels;
    ROC curve, choice of threshold; compare energy-based, FFT-based and
    MEM-based features.}

\exercise{12.3}{Real-time knock detector}{\W}{
  \item ESP32 + \texttt{crank-60-2} + \texttt{knock-sensor} chips: full
    chain in real time, knock LED per cylinder, ignition retard as
    output; vary RPM and knock intensity.
  \item Optional: port to Uno in fixed point and report the limits.}

%=====================================================================
\section*{Open points}
\begin{itemize}
  \item Confirm the JSLinux image and toolchain to be used
    (\texttt{gcc} vs.\ \texttt{tcc}; Python is not assumed).
  \item Confirm which Wokwi features students may use (custom chips,
    ESP32, project sharing) under the chosen account type.
  \item Fix the numerical parameters (knock frequencies, sampling rates)
    once Chapters 4, 5 and 12 of the script are available.
  \item Language of the sheets (English, as the script).
\end{itemize}

\end{document}
